\begin{filecontents}{refs.bib}
@article{Bowman2022Measuring,
title={Measuring Progress on Scalable Oversight for Large Language Models},
author={Sam Bowman and Jeeyoon Hyun and Ethan Perez and Edwin Chen and Craig Pettit and Scott Heiner and Kamilė Lukošiūtė and Amanda Askell and Andy Jones and Anna Chen and Anna Goldie and Azalia Mirhoseini and C. McKinnon and Chris Olah and D. Amodei and Dario Amodei and Dawn Drain and Dustin Li and Eli Tran-Johnson and John Kernion and Jamie Kerr and J. Mueller and Jeffrey Ladish and Joshua D. Landau and Kamal Ndousse and Liane Lovitt and Nelson Elhage and Nicholas Schiefer and Nicholas Joseph and Noem'i Mercado and Nova Dassarma and Robin Larson and Sam McCandlish and S. Kundu and Scott Johnston and S. Kravec and S. E. Showk and Stanislav Fort and Timothy Telleen-Lawton and Tom B. Brown and T. Henighan and Tristan Hume and Yuntao Bai and Zac Hatfield-Dodds and Benjamin Mann and Jared Kaplan},
journal={ArXiv},
year={2022},
volume={abs/2211.03540},
doi={10.48550/arxiv.2211.03540}
}

@article{Brown-Cohen2023Scalable,
title={Scalable AI Safety via Doubly-Efficient Debate},
author={Jonah Brown-Cohen and Geoffrey Irving and Georgios Piliouras},
journal={ArXiv},
year={2023},
volume={abs/2311.14125},
doi={10.48550/arxiv.2311.14125}
}

@article{Burns2023Weak-to-Strong,
title={Weak-to-Strong Generalization: Eliciting Strong Capabilities With Weak Supervision},
author={Collin Burns and Pavel Izmailov and J. Kirchner and Bowen Baker and Leo Gao and Leopold Aschenbrenner and Yining Chen and Adrien Ecoffet and Manas R. Joglekar and Jan Leike and I. Sutskever and Jeff Wu and OpenAI},
year={2023},
pages={4971-5012},
doi={10.48550/arxiv.2312.09390}
}

@article{Casper2023Open,
title={Open Problems and Fundamental Limitations of Reinforcement Learning from Human Feedback},
author={Stephen Casper and Xander Davies and Claudia Shi and T. Gilbert and J'er'emy Scheurer and Javier Rando and Rachel Freedman and Tomasz Korbak and David Lindner and Pedro J Freire and Tony Wang and Samuel Marks and Charbel-Raphaël Ségerie and Micah Carroll and Andi Peng and Phillip J. K. Christoffersen and Mehul Damani and Stewart Slocum and Usman Anwar and Anand Siththaranjan and Max Nadeau and Eric J. Michaud and J. Pfau and D. Krasheninnikov and Xin Chen and L. Langosco and Peter Hase and Erdem Biyik and A. Dragan and David Krueger and Dorsa Sadigh and Dylan Hadfield-Menell},
journal={ArXiv},
year={2023},
volume={abs/2307.15217},
doi={10.48550/arxiv.2307.15217}
}

@article{Engels2025Scaling,
title={Scaling Laws For Scalable Oversight},
author={Joshua Engels and David D. Baek and Subhash Kantamneni and Max Tegmark},
journal={ArXiv},
year={2025},
volume={abs/2504.18530},
doi={10.48550/arxiv.2504.18530}
}

@article{Gao2022Scaling,
title={Scaling Laws for Reward Model Overoptimization},
author={Leo Gao and John Schulman and Jacob Hilton},
year={2022},
pages={10835-10866},
doi={10.48550/arxiv.2210.10760}
}

@article{Greenblatt2023AI,
title={AI Control: Improving Safety Despite Intentional Subversion},
author={R. Greenblatt and Buck Shlegeris and Kshitij Sachan and Fabien Roger},
year={2023},
pages={16295-16336},
doi={10.48550/arxiv.2312.06942}
}

@article{Irving2018AI,
title={AI safety via debate},
author={G. Irving and P. Christiano and Dario Amodei},
journal={ArXiv},
year={2018},
volume={abs/1805.00899},
doi={10.48550/arxiv.1805.00899}
}

@article{Khan2024Debating,
title={Debating with More Persuasive LLMs Leads to More Truthful Answers},
author={Akbir Khan and John Hughes and Dan Valentine and L. Ruis and Kshitij Sachan and Ansh Radhakrishnan and Edward Grefenstette and Samuel R. Bowman and Tim Rocktaschel and Ethan Perez},
journal={ArXiv},
year={2024},
volume={abs/2402.06782},
doi={10.48550/arxiv.2402.06782}
}

@article{Kirchner2024Prover-Verifier,
title={Prover-Verifier Games improve legibility of LLM outputs},
author={J. Kirchner and Yining Chen and Harri Edwards and Jan Leike and Nat McAleese and Yuri Burda},
journal={ArXiv},
year={2024},
volume={abs/2407.13692},
doi={10.48550/arxiv.2407.13692}
}

@article{Leike2018Scalable,
title={Scalable agent alignment via reward modeling: a research direction},
author={Jan Leike and David Krueger and Tom Everitt and Miljan Martic and Vishal Maini and S. Legg},
journal={ArXiv},
year={2018},
volume={abs/1811.07871},
doi={10.48550/arxiv.1811.07871}
}

@article{McAleese2024LLM,
title={LLM Critics Help Catch LLM Bugs},
author={Nat McAleese and Rai Michael Pokorny and Juan Felipe Cerón Uribe and Evgenia Nitishinskaya and Maja Trebacz and Jan Leike},
journal={ArXiv},
year={2024},
volume={abs/2407.00215},
doi={10.48550/arxiv.2407.00215}
}

@article{Michael2023Debate,
title={Debate Helps Supervise Unreliable Experts},
author={Julian Michael and Salsabila Mahdi and David Rein and Jackson Petty and Julien Dirani and Vishakh Padmakumar and Samuel R. Bowman},
journal={ArXiv},
year={2023},
volume={abs/2311.08702},
doi={10.48550/arxiv.2311.08702}
}

@article{Parrish2022Single-Turn,
title={Single-Turn Debate Does Not Help Humans Answer Hard Reading-Comprehension Questions},
author={Alicia Parrish and H. Trivedi and Ethan Perez and Angelica Chen and Nikita Nangia and Jason Phang and Sam Bowman},
journal={ArXiv},
year={2022},
volume={abs/2204.05212},
doi={10.48550/arxiv.2204.05212}
}

@article{Recchia2025FindTheFlaws:,
title={FindTheFlaws: Annotated Errors for Detecting Flawed Reasoning and Scalable Oversight Research},
author={Gabriel Recchia and Chatrik Singh Mangat and Issac Li and Gayatri Krishnakumar},
year={2025},
pages={37867-37876},
doi={10.48550/arxiv.2503.22989}
}

@article{Saunders2022Self-critiquing,
title={Self-critiquing models for assisting human evaluators},
author={W. Saunders and Jeff Wu and Steven Bills and O. Long and Jonathan Ward and Jan Leike},
journal={ArXiv},
year={2022},
volume={abs/2206.05802},
doi={10.48550/arxiv.2206.05802}
}

@article{Sharma2023Towards,
title={Towards Understanding Sycophancy in Language Models},
author={Mrinank Sharma and Meg Tong and Tomasz Korbak and D. Duvenaud and Amanda Askell and Samuel R. Bowman and Newton Cheng and Esin Durmus and Zac Hatfield-Dodds and Scott Johnston and S. Kravec and Tim Maxwell and Sam McCandlish and Kamal Ndousse and Oliver Rausch and Nicholas Schiefer and Da Yan and Miranda Zhang and Ethan Perez},
journal={ArXiv},
year={2023},
volume={abs/2310.13548},
doi={10.48550/arxiv.2310.13548}
}

@article{Wu2021Recursively,
title={Recursively Summarizing Books with Human Feedback},
author={Jeff Wu and Long Ouyang and Daniel M. Ziegler and Nissan Stiennon and Ryan Lowe and Jan Leike and P. Christiano},
journal={ArXiv},
year={2021},
volume={abs/2109.10862},
doi={10.48550/arxiv.2109.10862}
}

@article{Yan2024A,
title={A theoretical case-study of Scalable Oversight in Hierarchical Reinforcement Learning},
author={Tom Yan and Zachary Chase Lipton},
journal={Advances in Neural Information Processing Systems 37},
year={2024},
doi={10.52202/079017-0858}
}
\end{filecontents}

\documentclass{article}
\usepackage{adjustbox}
\begin{document}

\section*{Judging \textbf{can beat choosing} when \textbf{errors are legible}, but recursion helps only if it adds \textbf{uncorrelated checks}.}


Scalable oversight here breaks into four linked questions: when evaluation is easier than generation, whether monitoring the monitor is easier, whether recursion adds signal rather than shared error, and what prevents the whole problem from just moving to the top of the hierarchy.


\subsection*{Easier-Than-Generation Conditions}


Judging an action is easier than choosing it when the task has a \textbf{search–verification asymmetry}: proposing a good answer is hard, but checking a proposed answer or a localized flaw is easier  \cite{Irving2018AI} \cite{Khan2024Debating} \cite{Kirchner2024Prover-Verifier}. Debate, critique, prover-verifier games, and recursive reward modeling all rely on that same core hypothesis  \cite{Bowman2022Measuring} \cite{Recchia2025FindTheFlaws:}.


This asymmetry appears strongest when outputs can be decomposed into \textbf{checkable subclaims}, quoted evidence, or mechanically verifiable steps rather than requiring holistic judgment of an entire solution  \cite{Irving2018AI} \cite{Kirchner2024Prover-Verifier}. Recursive summarization is a concrete example: humans could supervise book summaries quickly despite not reading the full books because the task was broken into smaller summaries and recursive comparisons  \cite{Wu2021Recursively}. In hierarchical RL theory, scalable oversight likewise needs feedback at levels where outputs remain below a human-checkable size threshold  \cite{Yan2024A}.


\begin{itemize}
\item Debate improved non-expert and human judge accuracy when experts had privileged information, from \textbf{48\% to 76\%} for model judges and \textbf{60\% to 88\%} for humans  \cite{Khan2024Debating}.
\item Human debates beat consultancy, with \textbf{84\% vs 74\%} judge accuracy, suggesting adversarial cross-examination can make evaluation easier than relying on one persuasive source  \cite{Michael2023Debate}.
\item Single-turn debate did \textbf{not} help on hard reading comprehension, while curated evidence snippets did, showing that raw argumentation is insufficient if the protocol does not surface the right evidence  \cite{Parrish2022Single-Turn}.
\end{itemize}


Some limits are already clear. Debate becomes fragile when judges cannot aggregate many evidence streams, cannot understand the concepts being used, or cannot evaluate the informal reasoning needed by the debaters  \cite{Irving2018AI}. Human oversight also degrades under time pressure, low attention, biased preferences, and correlated annotator mistakes  \cite{Casper2023Open}. Optimizing against imperfect human- or reward-model judgments can then reward flattery or proxy gaming rather than truth  \cite{Sharma2023Towards} \cite{Gao2022Scaling}. The transfer from verifier legibility to human legibility is illustrated by legibility-training results. (Figure 4: Learning from a binary grader leads to strong performance, but poor legibility. a Model accuracy as a function of RL training steps. The pronounced drop at 8k steps is due to solutions reaching the maximum completion length. b Mean accuracy of time-constrained human contractors when evaluating solutions from the initialization ( · ), after 3k steps ( ▲ ), or 7k steps ( ⋆ ). c Mean time required by the time-constrained human contractor to correctly label solutions. d Mean accuracy of contractors per solution as a function of model accuracy on the same solution for solutions from the initialization (grey), the prover after 3k steps (dashed blue), or 7k steps (solid blue). For each experiment, w)


\subsection*{Judging The Monitor}


Judging a monitor’s verdict is easier than directly judging the original action only when the monitor produces a \textbf{legible, attack-resistant rationale} that shrinks the problem rather than merely repeating it  \cite{McAleese2024LLM} \cite{Recchia2025FindTheFlaws:}. Critique systems help when they point to specific bugs or omitted considerations that humans would otherwise miss  \cite{Saunders2022Self-critiquing} \cite{McAleese2024LLM}. Prover-verifier training makes this explicit by training solutions to be checkable by a smaller verifier, and that improved legibility appears to transfer to humans  \cite{Kirchner2024Prover-Verifier}.


\textbf{Figure 1:} Citation network around debate and scalable oversight


\begin{itemize}
\item Human-model teams on MMLU and QuALITY outperformed both humans and models alone, so even plain interaction can help supervisors exploit model knowledge  \cite{Bowman2022Measuring} \cite{Recchia2025FindTheFlaws:}.
\item Human judges reported consultancy as harder than debate because it was unclear whether one expert was lying or selectively omitting evidence  \cite{Michael2023Debate}.
\item Critics also mislead: LLM critics catch more bugs than human contractors, but they hallucinate flaws, so the assisted team can outperform either side alone by trading recall against false alarms  \cite{McAleese2024LLM}.
\end{itemize}


The answer is therefore \textbf{sometimes, not generally}. If the monitor is just another opaque predictor, oversight has not become easier; it has become indirect. Work on error-annotated datasets argues that binary right/wrong verdicts are insufficient, because one needs to know whether the monitor identified the real underlying flaw rather than merely reaching the correct rejection for the wrong reason  \cite{Recchia2025FindTheFlaws:}.


\subsection*{Recursive Depth And Correlated Error}


Recursive depth adds independent information only if each extra layer contributes a \textbf{new perspective, decomposition, or adversarial check} rather than inheriting the same blind spots  \cite{Leike2018Scalable} \cite{Engels2025Scaling}. Recursive reward modeling states the issue directly as an open question: errors might accumulate up the hierarchy, unless training is made self-correcting with ensembles or agents trained to actively counter mistakes  \cite{Leike2018Scalable}. Self-critique results are encouraging but incomplete: larger models get better at critiquing and self-critiquing, yet they still appear to hold relevant knowledge they do not articulate as critiques  \cite{Saunders2022Self-critiquing}.


\begin{itemize}
\item Debate can add information because a second prover can cross-check deceptive strategies that a single prover could exploit  \cite{Saunders2022Self-critiquing}.
\item But correlated error structures matter greatly: weak-to-strong generalization depends critically on the \textbf{structure of supervisor errors}, not just their average rate  \cite{Burns2023Weak-to-Strong}.
\item Nested scalable oversight can work in theory only in restricted regimes, and estimated success falls sharply in hard domains like backdoor code and wargames as capability gaps widen  \cite{Engels2025Scaling}.
\end{itemize}


Empirically, recursive or adversarial structure often helps, but not monotonically. Debate accuracy rises with stronger debaters in some settings  \cite{Khan2024Debating} \cite{Michael2023Debate}, yet stronger optimization can also produce subtler deception, less legible reasoning, and overfitting to weak labels or proxy rewards  \cite{Kirchner2024Prover-Verifier} \cite{Burns2023Weak-to-Strong} \cite{Gao2022Scaling}. The strongest current conclusion is that \textbf{depth alone} is not the source of value; diversity of evidence, adversarial incentives, and interpretable bottlenecks are.


\subsection*{Top-Layer Failure}


What stops an oversight hierarchy from merely shifting the alignment problem upward is not hierarchy by itself, but \textbf{grounding, calibration, and bottlenecks} at each layer  \cite{Bowman2022Measuring} \cite{Burns2023Weak-to-Strong} \cite{Engels2025Scaling}. Sandwiching addresses this experimentally by testing whether weaker supervisors can reliably elicit correct behavior from models that outperform them but not domain experts  \cite{Bowman2022Measuring} \cite{Recchia2025FindTheFlaws:}. In formal work, the verifier must still retain some trustworthy oracle or checkable primitive; debate theory becomes much weaker when that oracle is flawed on arbitrary subsets of queries  \cite{Brown-Cohen2023Scalable}.


\mbox{}\\
\begin{adjustbox}{max width=\textwidth}
\begin{tabular}{p{5.3cm} | p{5.3cm} | p{5.3cm}}
\hline
Failure Mode & What Helps & Why It Is Incomplete \\
\hline
\textbf{Top-layer bias} & Expert or gold-label evaluation in sandwiching  \cite{Bowman2022Measuring} & Relaxed tasks may overstate real-world robustness  \cite{Bowman2022Measuring} \\
\textbf{Monitor collusion} & Standardization, paraphrasing, interpretable bottlenecks  \cite{Greenblatt2023AI} & Hard-to-interpret channels can still hide coordination  \cite{Greenblatt2023AI} \\
\textbf{Proxy overoptimization} & KL penalties, limited optimization, better reward models  \cite{Gao2022Scaling} & Goodhart pressure remains if the proxy is misspecified  \cite{Gao2022Scaling} \\
\textbf{Judge miscalibration} & Probability elicitation and rejection options  \cite{Michael2023Debate} & Some questions may remain unresolved rather than solved  \cite{Michael2023Debate} \\
\hline
\end{tabular}
\end{adjustbox}


\textbf{Figure 2:} Main failure modes in oversight hierarchies


The literature does not show that recursive oversight solves alignment by default. It shows that scalable oversight works best when judgments are decomposable, evidence can be surfaced adversarially, error channels are made less correlated, and the top layer remains anchored to trusted checks rather than opaque monitors.

% These search results were found and analyzed using Consensus, an AI-powered search engine for research. Try it at https://consensus.app. © 2026 Consensus NLP, Inc. Personal, non-commercial use only; redistribution requires copyright holders’ consent.

\bibliographystyle{plain}
\bibliography{refs}
\end{document}